\section{第 27 轮：秩一截面的 Fourier 闭合——无条件 $\mathrm{Neg}\le\dfrac5{6\pi}$ 与 $\kappa_3(\text{秩一})<\rho_0$}\label{sec:r27}

第 26 轮把秩一截面化成了一维量 $L_w=1+2\mathrm{Neg}$ 的上界问题，并在\emph{锥帧}上求出
$\sup L_w=\frac{13+\sqrt5}{10}$（推论~\ref{cor:r26rank1cone}），但「极值元必为锥」这条对称化命题
（注~\ref{rem:r26one}~的 P1）始终没有进展。本轮\textbf{绕开}了对称化的依赖：
我们直接在\emph{一般}秩一帧上，用一条「Young 不等式 $+$ 刚性恒等式 $+$ 圆上 Fourier」的链条，
证明
\begin{equation}\label{eq:r27head}
\boxed{\;\mathrm{Neg}\ \le\ \frac5{6\pi}=0.265258238486492\;<\;\rho_0-1=0.281568196872168\;},
\end{equation}
从而
\begin{equation}\label{eq:r27kappa}
\kappa_3(\text{秩一})=1+\mathrm{Neg}\ \le\ 1+\frac5{6\pi}=1.265258238486492\;<\;\rho_0 ,
\end{equation}
\textbf{无条件}地闭合了秩一截面（定理~\ref{thm:r27main}）。
这条结果有三点值得强调。
第一，它\emph{不需要}锥最优性：对一切满足各向同性的秩一帧都成立，
因此把第 26 轮「条件性」的 $\frac{11+\varphi}{10}=1.261803$~升级为「无条件」的 $1+\frac5{6\pi}=1.265258$，
二者只差 $3.45\times10^{-3}$，而后者是\textbf{已证}的。
第二，它把 \STW\ 的逐点天花板 $\Phi=1.309017$~在秩一截面上改进了 $4.38\times10^{-2}$，
这是本课题第一次在\emph{某个截面}上把上界压到 $\rho_0$ 以下并给出严格证明。
第三，证明的核心是一个\textbf{刚性恒等式}（引理~\ref{lem:r27rigid}）
$\sum_{ij}\rho_i\rho_jc_{ij}^2=\|I_2\|_F^2=2$
与一个\textbf{Fourier 巧合}：钝角余弦平方核 $K(u)=(\cos u)_-^2$~的第二系数 $\widehat K_2=\frac18$~
恰好被二维各向同性 $\sum\rho_ie^{2i\varphi_i}=0$~杀掉，而 $\widehat K_2$~正是除常数项外最大的正系数。
这一「各向同性自动消灭最坏频率」的现象，是本轮方法论上最重要的观察，
也是我们把 P1（对称化）暂时搁置、转而直接估计 $\mathrm{Neg}$~的理由。

本轮另附一条独立的引理~\ref{lem:r27zs}：$\|Z_S\|^2\le\sigma_S$~对一切子集 $S$~成立，
它给出秩一刚性 $A_w\equiv1$（命题~\ref{prop:r25rank1}）的一个\emph{一行}新证明，
只用 Cauchy--Schwarz 与 $\sum_S y_iy_i^\top\preceq I_3$，不依赖第 25 轮的偶极子机器。

\subsection{秩一截面的双权重坐标}\label{sec:r27:setup}

设 $\Lam=zz^\top$（$\|z\|=1$）。经引理~\ref{lem:r26hemi}~的半球翻转，可设
$a_i:=\langle u_i,z\rangle\ge0$、$b_i:=\sqrt{1-a_i^2}$，并把 $u_i$~在 $z^\perp$~上的单位化分量记为
$p_i\in S^1\subset z^\perp$，方位角 $\varphi_i$，于是
\[
u_i=a_iz+b_ip_i,\qquad
\langle u_i,u_j\rangle=a_ia_j+b_ib_jc_{ij},\qquad
c_{ij}:=\langle p_i,p_j\rangle=\cos(\varphi_i-\varphi_j).
\]
引入\textbf{双权重}（沿用注~\ref{rem:r26struct}(2)~的记号）
\begin{equation}\label{eq:r27var}
\sigma_i:=g_ia_i^2\ge0,\qquad
\rho_i:=g_ib_i^2\ge0,\qquad
r_i:=g_ia_ib_i=\sqrt{\sigma_i\rho_i},\qquad
\lambda_i:=\frac{a_i}{b_i}=\sqrt{\frac{\sigma_i}{\rho_i}} .
\end{equation}
帧条件 $\sum_ig_iu_iu_i^\top=I_3$、$\sum_ig_i=3$~沿 $z\oplus z^\perp$~分解为四条约束：
\begin{align}
&\text{(C1)}\quad \sum_i\sigma_i=\sum_ig_ia_i^2=1, \label{eq:r27C1}\\
&\text{(C2)}\quad \sum_i\rho_i=\sum_ig_ib_i^2=\sum_ig_i-\sum_i\sigma_i=2, \label{eq:r27C2}\\
&\text{(C3)}\quad \sum_i\rho_ip_ip_i^\top=I_2
\iff \sum_i\rho_ie^{2i\varphi_i}=0\ \ (\text{无迹部分}), \label{eq:r27C3}\\
&\text{(C4)}\quad \sum_ir_ip_i=\sum_ig_ia_ib_ip_i=0
\iff \sum_ir_ie^{i\varphi_i}=0\ \ (\text{交叉项}). \label{eq:r27C4}
\end{align}
这里 (C3) 是 $z^\perp$~块 $=$\,$I_2$，(C4) 是 $z$--$z^\perp$~交叉块 $=0$；
(C3) 的\textbf{无迹部分} $\sum\rho_ie^{2i\varphi_i}=0$~将在引理~\ref{lem:r27fourier}~里起决定作用。

命题~\ref{prop:r26rank1}~已给出 $L_w=1+2\mathrm{Neg}$，其中
\begin{equation}\label{eq:r27neg}
\mathrm{Neg}=\sum_{\langle u_i,u_j\rangle<0}g_ig_ja_ia_j\,|\langle u_i,u_j\rangle|
=\sum_{(i,j)\in\mathcal N}(-t_{ij}),\qquad
t_{ij}:=w_iw_j\langle u_i,u_j\rangle=\sigma_i\sigma_j+r_ir_jc_{ij},
\end{equation}
$w_i=g_ia_i$，负集
\begin{equation}\label{eq:r27Nset}
\mathcal N:=\{(i,j):t_{ij}<0\}=\{(i,j):c_{ij}<-\lambda_i\lambda_j\}\subseteq\{(i,j):c_{ij}<0\}.
\end{equation}
（最后一步用 $t_{ij}<0\iff r_ir_jc_{ij}<-\sigma_i\sigma_j\iff c_{ij}<-\frac{\sigma_i\sigma_j}{r_ir_j}=-\lambda_i\lambda_j$。）
两条平凡的恒等式贯穿全轮：
\begin{equation}\label{eq:r27trivial}
\sum_{ij}t_{ij}=\Bigl(\sum_i\sigma_i\Bigr)^2+\Bigl|\sum_ir_ie^{i\varphi_i}\Bigr|^2\cdot(\text{实部})=1
\quad(\text{由 (C1)(C4)}),\qquad
L_w=\sum_{ij}|t_{ij}|=1+2\mathrm{Neg}.
\end{equation}

\subsection{Young 不等式与 $\kappa-2$ 亏格恒等式}\label{sec:r27:young}

\begin{lemma}[Young 逐点估计与精确亏格]\label{lem:r27young}
对 $(i,j)\in\mathcal N$，记 $A:=r_ir_j|c_{ij}|\ge0$、$B:=\sigma_i\sigma_j>0$、
$\kappa_{ij}:=A/B=\dfrac{|c_{ij}|}{\lambda_i\lambda_j}$。则 $-t_{ij}=A-B$，且
\begin{equation}\label{eq:r27young}
-t_{ij}=A-B\le\frac{A^2}{4B}=\frac{\rho_i\rho_jc_{ij}^2}{4},
\qquad\text{等号}\iff\kappa_{ij}=2,
\end{equation}
更精确地有\textbf{亏格恒等式}
\begin{equation}\label{eq:r27defect}
\mathrm{Neg}=\frac\Theta4-\frac14\sum_{(i,j)\in\mathcal N}\sigma_i\sigma_j(\kappa_{ij}-2)^2,
\qquad
\Theta:=\sum_{(i,j)\in\mathcal N}\rho_i\rho_jc_{ij}^2 .
\end{equation}
\end{lemma}

\begin{proof}
$c_{ij}<0$~于 $\mathcal N$，故 $-t_{ij}=-(\sigma_i\sigma_j+r_ir_jc_{ij})=r_ir_j|c_{ij}|-\sigma_i\sigma_j=A-B$。
初等恒等式 $A-B=\frac{A^2}{4B}-\frac{(A-2B)^2}{4B}$~（两边同乘 $4B$~即
$4AB-4B^2=A^2-(A^2-4AB+4B^2)$）给出 \eqref{eq:r27young}~与等号条件 $A=2B\iff\kappa_{ij}=2$。
又 $\frac{A^2}{4B}=\frac{r_i^2r_j^2c_{ij}^2}{4\sigma_i\sigma_j}
=\frac{(\sigma_i\rho_i)(\sigma_j\rho_j)c_{ij}^2}{4\sigma_i\sigma_j}=\frac{\rho_i\rho_jc_{ij}^2}4$，
且 $\frac{(A-2B)^2}{4B}=\frac{B^2(\kappa_{ij}-2)^2}{4B}=\frac{\sigma_i\sigma_j(\kappa_{ij}-2)^2}4$。
对 $(i,j)\in\mathcal N$~求和即得 \eqref{eq:r27defect}。
\end{proof}

\begin{corollary}[$\mathrm{Neg}\le P/4$]\label{cor:r27P}
由 $\mathcal N\subseteq\{c_{ij}<0\}$~且被求和项非负，
\begin{equation}\label{eq:r27Pdef}
\Theta\le P:=\!\!\sum_{\substack{i\ne j\\ c_{ij}<0}}\!\!\rho_i\rho_jc_{ij}^2,
\qquad\text{故}\qquad
\mathrm{Neg}\le\frac P4 .
\end{equation}
（对角项 $c_{ii}=1>0$~不进入 $P$，故 $i\ne j$~的限制无影响。）
等号 $\Theta=P\iff\mathcal N=\{c_{ij}<0\}\iff$~每个钝角对都满足 $c_{ij}<-\lambda_i\lambda_j$；
冠军处恰好取等（注~\ref{rem:r26struct}(1)：负对满足 $\langle u_i,u_j\rangle=-a_ia_j$）。
\end{corollary}

\subsection{刚性恒等式}\label{sec:r27:rigid}

\begin{lemma}[刚性恒等式]\label{lem:r27rigid}
在约束 (C3) 下，
\begin{equation}\label{eq:r27rigid}
\sum_{i,j}\rho_i\rho_jc_{ij}^2=\Bigl\|\sum_i\rho_ip_ip_i^\top\Bigr\|_F^2=\|I_2\|_F^2=2 .
\end{equation}
特别地，$P\le2$~（粗界），且全和是\textbf{常数}（不依赖帧），只有「落在钝角对上的部分」$P$~是变量。
\end{lemma}

\begin{proof}
$c_{ij}^2=\langle p_i,p_j\rangle^2=\tr(p_ip_i^\top p_jp_j^\top)$，故
$\sum_{ij}\rho_i\rho_jc_{ij}^2=\tr\bigl[(\sum_i\rho_ip_ip_i^\top)^2\bigr]
=\|\sum_i\rho_ip_ip_i^\top\|_F^2$（对称阵）。由 (C3) 该矩阵 $=I_2$，$\|I_2\|_F^2=2$。
$P$~是把全和中 $c_{ij}\ge0$~的项丢掉，故 $P\le2$。
\end{proof}

粗界 $P\le2$~只给 $\mathrm{Neg}\le\frac12=0.5$，远不足以关闭（目标 $\rho_0-1=0.2816$）。
关键改进来自下一节的 Fourier 分析：它把「$c_{ij}\ge0$~的项被丢掉」这件事\emph{定量化}，
并利用 (C3) 的无迹部分。

\subsection{$P$ 的 Fourier 上界}\label{sec:r27:fourier}

把 $\mu:=\sum_i\rho_i\delta_{\varphi_i}$~看作 $S^1$~上的正测度，$\mu(S^1)=\sum\rho_i=2$（C2），
$\widehat\mu(k):=\int e^{ik\varphi}\dd\mu=\sum_i\rho_ie^{ik\varphi_i}$，
于是 $\widehat\mu(0)=2$、$\widehat\mu(2)=0$（C3 的无迹部分）。
定义核
\begin{equation}\label{eq:r27K}
K(u):=(\cos u)_-^2=\begin{cases}\cos^2u,&\cos u<0,\\[2pt]0,&\cos u\ge0,\end{cases}
\qquad
P=\iint_{S^1\times S^1}K(\varphi-\varphi')\,\dd\mu(\varphi)\dd\mu(\varphi'),
\end{equation}
（因为 $K(\varphi_i-\varphi_j)=c_{ij}^2\cdot1\{c_{ij}<0\}$，而对角项 $K(0)=0$）。

\begin{lemma}[$K$~的 Fourier 系数]\label{lem:r27Khat}
$K$~是实偶函数，$K(u)=\sum_{k\in\Z}\widehat K_ke^{iku}$，$\widehat K_{-k}=\widehat K_k\in\R$，且
\begin{equation}\label{eq:r27Khat}
\widehat K_0=\frac14,\qquad
\widehat K_2=\frac18,\qquad
\widehat K_k=0\ (k\ \text{偶},\,k\ge4),\qquad
\widehat K_k=\frac{2\sin(k\pi/2)}{\pi k(k^2-4)}\ (k\ \text{奇}).
\end{equation}
特别地，对奇数 $k$：$\widehat K_k>0\iff k\equiv1\ (\mathrm{mod}\ 4),\,k\ge5$；
$\widehat K_1=-\frac2{3\pi}<0$；$\widehat K_k<0\iff k\equiv3\ (\mathrm{mod}\ 4)$。
\end{lemma}

\begin{proof}
由 $\cos^2u=\frac{1+\cos2u}2=\frac12+\frac{e^{2iu}+e^{-2iu}}4$~与
$\int_{\pi/2}^{3\pi/2}e^{imu}\dd u=e^{im\pi/2}\frac{e^{im\pi}-1}{im}$~（$m\ne0$）、$=pi$~当 $m=0$，
在 $K$~的支撑 $[\frac\pi2,\frac{3\pi}2]$~上逐项积分即得。
$k$~偶 $\ge4$~时三个频率 $-k,2-k,-2-k$~都是非零偶数，$e^{im\pi}-1=0$，故 $\widehat K_k=0$；
$k=2$~时出现 $m=0$~项，给 $\widehat K_2=\frac1{2\pi}\cdot\frac\pi4=\frac18$；
$k$~奇时整理得 $\widehat K_k=\frac{2\sin(k\pi/2)}{\pi k(k^2-4)}$。
证书 D 段以 $2\times10^6$~点数值积分核对 $k=0,\dots,9$，最大差 $1.4\times10^{-17}$。
\end{proof}

\begin{lemma}[$P$~的 Fourier 表示与上界]\label{lem:r27fourier}
\begin{equation}\label{eq:r27Pfourier}
P=\sum_{k\in\Z}\widehat K_k\,|\widehat\mu(k)|^2
=1+2\!\!\sum_{\substack{k\ge1\\ k\ne2}}\!\!\widehat K_k|\widehat\mu(k)|^2 ,
\end{equation}
且
\begin{equation}\label{eq:r27Pbound}
P\ \le\ 1+8\!\!\sum_{\substack{k\ge1\ \text{奇}\\ \widehat K_k>0}}\!\!\widehat K_k
\ =\ 1+\frac{16}\pi S,\qquad
S:=\sum_{m\ge1}\frac1{(4m+1)\bigl((4m+1)^2-4\bigr)} .
\end{equation}
\end{lemma}

\begin{proof}
把 $K(u)=\sum_k\widehat K_ke^{iku}$~代入 \eqref{eq:r27K}~并交换积分与求和（$K$~连续、$\mu$~有限，绝对收敛）：
$P=\sum_k\widehat K_k\int\!\!\int e^{ik(\varphi-\varphi')}\dd\mu\dd\mu'
=\sum_k\widehat K_k|\widehat\mu(k)|^2$。
常数项 $\widehat K_0|\widehat\mu(0)|^2=\frac14\cdot4=1$；
$k=\pm2$~项 $\widehat K_2|\widehat\mu(2)|^2=\frac18\cdot0=0$——\textbf{这正是各向同性 (C3) 的作用}：
它自动消灭了除常数项外最大的正系数 $\widehat K_2=\frac18$。
其余项用 $|\widehat\mu(k)|\le\widehat\mu(0)=2$~与「丢掉 $\widehat K_k<0$~的项」放大：
$P\le1+2\sum_{k\ge1,k\ne2}\widehat K_k\cdot4\cdot1\{\widehat K_k>0\}
=1+8\sum_{k\ge1,\widehat K_k>0}\widehat K_k$
（$k=2$~虽 $\widehat K_2>0$~但已被 $|\widehat\mu(2)|=0$~杀掉，不在求和内）。
由引理~\ref{lem:r27Khat}，正的 $\widehat K_k$~恰为奇数 $k\equiv1\ (4)$、$k\ge5$，
此时 $\sin(k\pi/2)=1$、$\widehat K_k=\frac2{\pi k(k^2-4)}$，令 $k=4m+1$~即得 $S$。
\end{proof}

\begin{lemma}[$S$~的闭式]\label{lem:r27S}
\begin{equation}\label{eq:r27S}
S=\sum_{m\ge1}\frac1{(4m+1)((4m+1)^2-4)}=\frac5{24}-\frac\pi{16}=0.011983792483971255\ldots
\end{equation}
\end{lemma}

\begin{proof}
部分分式 $\frac1{(4m-1)(4m+1)(4m+5)}=-\frac1{8(4m+1)}+\frac1{12(4m-1)}+\frac1{24(4m+5)}$
把 $S$~化为三个调和型和，用 digamma 函数 $\psi$~求值：
$S=\frac1{16}\bigl[\psi(\tfrac14)-\psi(\tfrac34)\bigr]+\frac5{24}$。
由 Gauss 的 digamma 反射值 $\psi(\tfrac14)-\psi(\tfrac34)=-\pi$（等价于 $\psi(1-x)-\psi(x)=\pi\cot\pi x$~在 $x=\frac14$），
得 $S=\frac5{24}-\frac\pi{16}$。证书 A 段以 $60$~位 \texttt{mpmath}~核对，与部分和（$m<2\times10^5$）差 $2\times10^{-13}$。
\end{proof}

\begin{theorem}[$P$~的闭合上界]\label{thm:r27P}
\begin{equation}\label{eq:r27Pclosed}
P\ \le\ 1+\frac{16}\pi\Bigl(\frac5{24}-\frac\pi{16}\Bigr)=\frac{10}{3\pi}=1.061032953945969\ldots
\end{equation}
\end{theorem}

\begin{proof}
$1+\frac{16}\pi S=1+\frac{16}\pi(\frac5{24}-\frac\pi{16})=1+\frac{16\cdot5}{24\pi}-1=\frac{80}{24\pi}=\frac{10}{3\pi}$。
\end{proof}

\subsection{主定理：秩一截面的无条件闭合}\label{sec:r27:main}

\begin{theorem}[秩一截面的无条件上界，本轮主定理]\label{thm:r27main}
对\textbf{一切}满足各向同性的秩一帧 $(g,U)$（$\Lam=zz^\top$），
\begin{equation}\label{eq:r27main}
\mathrm{Neg}\ \le\ \frac P4\ \le\ \frac14\cdot\frac{10}{3\pi}=\frac5{6\pi}=0.265258238486492,
\end{equation}
从而
\begin{equation}\label{eq:r27concl}
\kappa_3(\text{秩一})=1+\mathrm{Neg}\ \le\ 1+\frac5{6\pi}=1.265258238486492\ <\ \rho_0=1.281568196872168,
\end{equation}
余量 $\rho_0-1-\frac5{6\pi}=0.016309958385676$。
等价地 $L_w\le1+\frac5{3\pi}=1.530516476972984$。
\end{theorem}

\begin{proof}
链条 $\mathrm{Neg}\le\frac\Theta4\le\frac P4$（引理~\ref{lem:r27young}、推论~\ref{cor:r27P}）
与 $P\le\frac{10}{3\pi}$（定理~\ref{thm:r27P}）给出 \eqref{eq:r27main}；
$\kappa_3(\text{秩一})=1+\mathrm{Neg}$~是命题~\ref{prop:r26rank1}。
\end{proof}

\begin{remark}[余量的严格性]\label{rem:r27margin}
$\frac5{6\pi}<\rho_0-1\iff\pi>\frac5{6(\rho_0-1)}=2.9596145537$，
而 $\pi>3.14$~是初等的（阿基米德）。故余量是\textbf{严格可证}的，不依赖任何数值近似。
更精细地，$\rho_0-1-\frac5{6\pi}=\frac{-15+\pi(-9+\sqrt5+2\sqrt{35})}{18\pi}$（$60$~位核对）。
\end{remark}

\begin{remark}[与 \STW\ 及锥族条件值的对照]\label{rem:r27compare}
\begin{center}\small
\begin{tabular}{lccc}
\toprule
秩一截面上界 & 数值 & 状态 & 来源\\
\midrule
\STW\ 逐点天花板 $\Phi=\frac{3+\sqrt5}4$ & $1.309017$ & 已证 & 定理~\ref{thm:zonal}\\
本轮 $1+\frac5{6\pi}$ & $1.265258$ & \textbf{已证（无条件）} & 定理~\ref{thm:r27main}\\
锥族条件值 $\frac{11+\varphi}{10}$ & $1.261803$ & 条件性（需 P1） & 推论~\ref{cor:r26rank1cone}\\
目标 $\rho_0$ & $1.281568$ & —— & ——\\
\bottomrule
\end{tabular}
\end{center}
本轮把秩一截面从「条件性 $1.261803$」变成「无条件 $1.265258$」，
代价是 $3.45\times10^{-3}$~的松弛，收益是\textbf{不再依赖锥最优性}。
相对 \STW\ 改进 $4.38\times10^{-2}$，相对目标 $\rho_0$~留有 $1.63\times10^{-2}$~余量。
\end{remark}

\subsection{秩一刚性的新证明：$\|Z_S\|^2\le\sigma_S$}\label{sec:r27:zs}

命题~\ref{prop:r25rank1}~（秩一截面 $A_w\equiv1$）原本用第 25 轮的偶极子机器证明。
本轮发现一个只依赖帧理论的\emph{一行}证明，记录如下。

\begin{lemma}[子集质量不等式]\label{lem:r27zs}
设 $z_1,\dots,z_N\in\R^3$~满足 $\sum_iz_i=e$（某单位向量）与 $\sum_iz_iz_i^\top/\sigma_i=I_3$，
其中 $\sigma_i>0$、$\sigma_S:=\sum_{i\in S}\sigma_i$、$Z_S:=\sum_{i\in S}z_i$。则对\textbf{一切}子集 $S$，
\begin{equation}\label{eq:r27zs}
\|Z_S\|^2\le\sigma_S .
\end{equation}
\end{lemma}

\begin{proof}
令 $y_i:=z_i/\sqrt{\sigma_i}$，则 $\sum_{i\in S}y_iy_i^\top=\sum_{i\in S}z_iz_i^\top/\sigma_i\preceq I_3$。
其 Gram 阵 $G_S=(\langle y_i,y_j\rangle)_{i,j\in S}$~与 $\sum_{i\in S}y_iy_i^\top$~有相同的非零特征值，
故 $G_S\preceq I_{|S|}$。而 $Z_S=\sum_{i\in S}\sqrt{\sigma_i}\,y_i$，于是
$\|Z_S\|^2=(\sqrt{\sigma})^\top G_S(\sqrt{\sigma})\le\|\sqrt\sigma\|^2=\sigma_S$。
\end{proof}

\begin{proposition}[$A_w=1$~的新证明]\label{prop:r27aw}
在秩一截面、半球翻转后，取 $z_i:=w_iu_i=g_ia_iu_i$（即 $z_i=v_{i,1}v_i$，$v_i$~为 Stiefel 行）。
则 $\sum_iz_i=z$、$\sum_iz_iz_i^\top/\sigma_i=I_3$（$\sigma_i=g_ia_i^2$），
且对任意符号 $\varepsilon\in\{\pm1\}^N$、$S=\{i:\varepsilon_i=+1\}$，
\[
\Bigl\|\sum_i\varepsilon_iz_i\Bigr\|^2=\|2Z_S-e\|^2=4\|Z_S\|^2-4\langle Z_S,e\rangle+1
=4\|Z_S\|^2-4\sigma_S+1\le1,
\]
（用 $\langle Z_S,e\rangle=\sum_{i\in S}w_i\langle u_i,z\rangle=\sum_{i\in S}g_ia_i^2=\sigma_S$~与引理~\ref{lem:r27zs}）。
全 $+$~号给 $\|\sum z_i\|^2=\|z\|^2=1$，故 $A_w=\max_\varepsilon\|\sum\varepsilon_iz_i\|=1$。
\end{proposition}

\begin{proof}
$\sum_iz_i=\sum_ig_ia_iu_i=\sum_ig_i\langle u_i,z\rangle u_i=(\sum_ig_iu_iu_i^\top)z=z$（引理~\ref{lem:r26hemi}(ii)）；
$\sum_iz_iz_i^\top/\sigma_i=\sum_i(g_ia_i)^2u_iu_i^\top/(g_ia_i^2)=\sum_ig_iu_iu_i^\top=I_3$。
其余见正文。证书 C 段在 $N=7$、$120$~个随机 Stiefel 帧、全部 $2^7-1$~个子集上核验
$\|Z_S\|^2\le\sigma_S$，违例 $0$、最大 margin $8.9\times10^{-16}$。
\end{proof}

\subsection{紧性讨论：冠军与 $P$~的真实上确界}\label{sec:r27:tight}

冠军 $W_6=K_1\vee C_5$~处（注~\ref{rem:r26struct}）：
\begin{itemize}\itemsep=1pt
\item Young 取等：$\kappa_{ij}=2$~对一切 $(i,j)\in\mathcal N$（亏格项 $=7.7\times10^{-33}$）；
\item $\mathcal N=\{c_{ij}<0\}$，故 $\Theta=P$；
\item $P_{\text{champ}}=\frac{3+\sqrt5}5=1.047213595499958$，$\mathrm{Neg}_{\text{champ}}=\frac{\varphi^2}{10}=0.261803398874989$。
\end{itemize}
因此链条 $\mathrm{Neg}=\frac\Theta4=\frac P4\le\frac14\cdot\frac{10}{3\pi}$~中，
\textbf{唯一的松弛}在 $P_{\text{champ}}=\frac{3+\sqrt5}5=1.047214<\frac{10}{3\pi}=1.061033$（差 $1.38\times10^{-2}$）。
这给出两条后续路线：
\begin{enumerate}\itemsep=1pt
\item[\textbf{(i)}]\textbf{把 Fourier 界收紧到冠军值}。$P=\sum_k\widehat K_k|\widehat\mu(k)|^2$~的上确界
是一个\emph{Toeplitz 矩问题}：在 $\mu\ge0$、$\widehat\mu(0)=2$、$\widehat\mu(2)=0$、
且 $(\widehat\mu(k))$~正定（Toeplitz 阵半正定）的约束下最大化 $P$。
冠军是 $5$-环 $\widehat\mu(k)=2\cdot1\{5\mid k\}$，它\emph{饱和}了 $|\widehat\mu(5)|=2$（$\widehat K_5=\frac2{105\pi}$~是最大正频率），
而本轮的界 $10/(3\pi)$~把\emph{所有} $k\equiv1(4)$~都按 $|\widehat\mu(k)|=2$~计，明显过松。
若能证明 $\sup P=\frac{3+\sqrt5}5$（$5$-环最优），则 $\mathrm{Neg}\le\frac{\varphi^2}{10}$、
$\kappa_3(\text{秩一})=\frac{11+\varphi}{10}$~\textbf{无条件且精确}。这是一个独立、可发表的目标（第 28 轮 P1$'$）。
\item[\textbf{(ii)}]\textbf{向单轴内部「增稠」}。见注~\ref{rem:r27global}。
\end{enumerate}

\begin{remark}[为什么本轮方法\emph{不能}直接推广到一般 $\Lam$]\label{rem:r27global}
Young 步 \eqref{eq:r27young}~的成立依赖 $t_{ij}$~有一个\textbf{正的} $\sigma_i\sigma_j$~偏移，
而该偏移来自秩一特有的 $h_i=a_i$（即 $w_i=g_ia_i$、$\sigma_i=g_ia_i^2=w_ia_i$）。
对一般单轴 $\Lam=bI_3+\rho ee^\top$，$h_i=\sqrt{b+\rho a_i^2}\ne a_i$，
即便半球翻转使 $a_i\ge0$、令 $\widetilde\sigma_i=w_ia_i\ge0$、$\widetilde r_i=w_ib_i$，
得到 $t_{ij}=\widetilde\sigma_i\widetilde\sigma_j+\widetilde r_i\widetilde r_jc_{ij}$，
Young 余项变为 $\frac{\widetilde r_i^2}{\widetilde\sigma_i}\cdot\frac{\widetilde r_j^2}{\widetilde\sigma_j}\cdot\frac{c_{ij}^2}4
=\bigl(g_ih_ib_i^2/a_i\bigr)\bigl(g_jh_jb_j^2/a_j\bigr)\frac{c_{ij}^2}4$，
其中权重 $g_ih_ib_i^2/a_i\ne\rho_i=g_ib_i^2$（除非 $h_i=a_i$）。
于是\textbf{刚性恒等式} $\sum\rho_i\rho_jc_{ij}^2=2$~不再适用于这个加权版本，
Fourier 步也就失去了「全和 $=$\,常数」这一支点。
更根本地，一般 $\Lam$~的最优符号模式 $\varepsilon^*$~是 $z_*$~的半球符号，
而 $z_*$~未必等于 $\Lam$~的轴 $e$；只有 $z_*=e$~时偏移 $\widetilde\sigma_i\widetilde\sigma_j$~才保持正号。
$z_*=e$~正是锥最优性（P1）的内容。
因此本轮的 Fourier 机器是\textbf{秩一专属}的；全局闭合仍须另寻支点（见第十四版队列）。
\end{remark}

\subsection{勘误}\label{sec:r27:erratum}

注~\ref{rem:r26struct}(1)~末句的公式应更正为
「每对的贡献 $g_1^2a^2\,(-\langle u_k,u_l\rangle)=g_1^2s_0^2$」——
原文多写了 $-a^2$（即 $g_1^2a^2(-\langle\cdot,\cdot\rangle-a^2)$），
代入 $-\langle u_k,u_l\rangle=s_0=a^2$~会使该式 $=0$，与 $=g_1^2s_0^2$~矛盾，系排版残留。
又「达到 $A=2B$~的比例」一句中的 $A=2B$~即引理~\ref{lem:r27young}~的 Young 取等条件
$\kappa_{ij}=A/B=2$（$A=r_ir_j|c_{ij}|$、$B=\sigma_i\sigma_j$），此处补注以免歧义。

\subsection{本轮状态、诚实标签与方法优先队列（第十四版）}\label{sec:r27:status}

\begin{center}\small
\begin{tabular}{lll}
\toprule
结论 & 状态 & 位置\\
\midrule
双权重坐标与 (C1)--(C4) & \textbf{已证（精确）} & \S\ref{sec:r27:setup}\\
Young 逐点 $+$ $\kappa-2$ 亏格恒等式 \eqref{eq:r27defect} & \textbf{已证（精确）} & 引理~\ref{lem:r27young}\\
$\mathrm{Neg}\le\Theta/4\le P/4$ & \textbf{已证} & 推论~\ref{cor:r27P}\\
刚性恒等式 $\sum\rho_i\rho_jc_{ij}^2=2$ & \textbf{已证（精确）} & 引理~\ref{lem:r27rigid}\\
$K$~的 Fourier 系数 \eqref{eq:r27Khat} & \textbf{已证（数值核对 $10^{-17}$）} & 引理~\ref{lem:r27Khat}\\
$P=\sum_k\widehat K_k|\widehat\mu(k)|^2$、各向同性杀 $\widehat K_2$ & \textbf{已证（精确）} & 引理~\ref{lem:r27fourier}\\
$S=\frac5{24}-\frac\pi{16}$ & \textbf{已证（符号 $60$ 位）} & 引理~\ref{lem:r27S}\\
$P\le\frac{10}{3\pi}$ & \textbf{已证（精确代数）} & 定理~\ref{thm:r27P}\\
$\mathrm{Neg}\le\frac5{6\pi}$、$\kappa_3(\text{秩一})\le1+\frac5{6\pi}<\rho_0$ & \textbf{已证（无条件）} & 定理~\ref{thm:r27main}\\
余量严格性（只需 $\pi>2.9597$） & \textbf{已证} & 注~\ref{rem:r27margin}\\
$\|Z_S\|^2\le\sigma_S\Rightarrow A_w=1$ & \textbf{已证（新证明）} & 引理~\ref{lem:r27zs}、命题~\ref{prop:r27aw}\\
$\sup P=\frac{3+\sqrt5}5$（$5$-环最优） & \textbf{猜想}（数值：冠军 $=$，随机搜索未超） & \S\ref{sec:r27:tight}(i)\\
\textbf{单轴内部／全局 $\kappa_3\le\rho_0$} & \textbf{未证}（秩一专属机器不可推广） & 注~\ref{rem:r27global}\\
\bottomrule
\end{tabular}
\end{center}

\textbf{方法优先队列（第十四版，数学优先；队首连续两次无进展则查文献并入队移植候选）}：
\begin{enumerate}\itemsep=2pt
\item[\textbf{P1$'$}]\textbf{（新，秩一精确值）证明 $\sup P=\frac{3+\sqrt5}5$}。
把 $P=\sum_k\widehat K_k|\widehat\mu(k)|^2$~的上确界写成 Toeplitz 矩问题，
用 LP/矩对偶找一张\emph{对偶证书}：即构造非负 Toeplitz 形式
$\sum_k\widehat K_k|\widehat\mu(k)|^2\le\frac{3+\sqrt5}5$~当 $\widehat\mu(0)=2,\widehat\mu(2)=0$。
入口：(a) 把约束 $(\widehat\mu(0),\widehat\mu(2))=(2,0)$~与正定性合并，对 $5$-环做互补松弛；
(b) Fej\'er--Riesz：$5$-环对应 $\mu=\frac25\sum_{k=0}^4\delta_{2\pi k/5}$，其湮灭多项式 $z^5-1$，
尝试证明任何可行 $\mu$~的 $P$~都被 $5$-环支配。
\textbf{硬校验}：对偶证书必须在 $5$-环取等。一旦成立，$\kappa_3(\text{秩一})=\frac{11+\varphi}{10}$~精确无条件。
\item[\textbf{P1}]\textbf{（沿用，全局关键路径）锥最优性}。
注~\ref{rem:r27global}~说明：全局闭合等价于「单轴截面极值元满足 $z_*=e$（即锥）」。
本轮的 Fourier 机器\emph{不能}替代它（刚性恒等式只在 $h_i=a_i$~时成立）。
入口沿用第十三版 (1a)--(1d)，但\textbf{优先级下调}：因为秩一截面已被本轮无条件闭合，
锥最优性现在只承担「单轴内部 $\rho\in(0,1)$」这一段，而不再是秩一与全局的公共缺口。
\item[\textbf{P2}]\textbf{（增稠）把秩一界推广到 $\rho$~接近 $1$~的单轴薄层}。
注~\ref{rem:r27global}~给出障碍因子 $h_i/a_i=\sqrt{\rho+(1-\rho)/(3a_i^2)}$。
若能在负集 $\mathcal N$~上证明 $h_i/a_i\le1+\varepsilon$（$\varepsilon\le0.0615$，由余量 $1.63\times10^{-2}$~反推），
则 $\mathrm{Neg}^*\le(1+\varepsilon)\frac5{6\pi}\le\rho_0-1$，单轴薄层闭合。
入口：负集限制 $c_{ij}<-\lambda_i\lambda_j$~与 $a_i$~的下界耦合（$a_i$~小则 $\lambda_i$~小、难进负集）。
\item[\textbf{P3}]\textbf{（窗口两尺度）}沿用第十二、十三版：$\mathrm{Neg}^*\le(\rho_0-1)+c\sin^2\theta$~中 $c=1$~的最优性。
\item[\textbf{P4}]\textbf{（中心余量）}沿用：$m_1=1.040\times10^{-3}$。
\item[\textbf{副}]\textbf{下界}：$n=4$~等纬度环给 $\rho=\frac54$~精确构造点；轮族极值性。
\end{enumerate}

\textbf{已弹出（累积，禁止再试）}：第九至十三版全部已弹出项，另加本轮一条：
\textbf{用刚性恒等式 $\sum\rho\rho c^2=2$~直接推广到一般 $\Lam$}——
该恒等式依赖 $h_i=a_i$（秩一专属），一般单轴的 Young 余项权重 $g_ih_ib_i^2/a_i\ne\rho_i$，
无对应刚性恒等式（注~\ref{rem:r27global}）。

复现：报告目录下 \texttt{证书\_秩一截面的Fourier闭合.py}（$32$~项全绿，输出 \texttt{\_cert\_r27\_out.txt}，约 $7$~秒）：
A 段 $60$~位常数代数（$S=\frac5{24}-\frac\pi{16}$、$1+\frac{16}\pi S=\frac{10}{3\pi}$、
$\mathrm{Neg}\le\frac5{6\pi}$、$\kappa_3(\text{秩一})\le1+\frac5{6\pi}<\rho_0$、余量闭式与严格性、与 \STW／锥族对照）；
B 段冠军精确恒等式（$\sigma/\rho/r/\lambda/t$、$\mathrm{Neg}=\frac{\varphi^2}{10}$、$L_w=\frac{13+\sqrt5}{10}$、
刚性恒等式 $=2$、$\Theta=P$、Young 取等 $\kappa=2$、$P_{\text{champ}}=\frac{3+\sqrt5}5$）；
C 段 $4000$~随机 Stiefel 帧（$\kappa-2$~亏格恒等式、刚性恒等式、$\mathrm{Neg}\le P/4$）
与 $\|Z_S\|^2\le\sigma_S$（$N=7$、$120$~帧、$2^7-1$~子集）；
D 段 Fourier 系数解析式 vs $2\times10^6$~点数值积分、$P=\sum_k\widehat K_k|\widehat\mu(k)|^2$~恒等式。
仓库根目录 \texttt{\_r27a.py}（冠军 $+$ 随机恒等式 $+$ $\|Z_S\|$）、
\texttt{\_r27b.py}（Fourier 系数 $+$ $S$~闭式 $+$ 最终界）、
\texttt{\_r27c.py}（符号 $S$、$P$~界塌缩、$\sup P$~数值探索）；
输出 \texttt{\_r27a\_out.txt}、\texttt{\_r27b\_out.txt}、\texttt{\_r27c\_out.txt}。

